\subsection{A full profinite lifting criterion and nonnilpotent finite tests}

\begin{claim}{D.23}{Proved}{Order-preserving lifting and conjugacy in the full shadow system}
Let $\mathcal I$ be the cofiltered poset of isolated finite-index normal
subgroups $N\unlhd B_3$ contained in the pure braid group $PB_3$, and put
$Q_N=GT(N)$.  The finite groups and their reduction homomorphisms
$R_{N,H}:Q_N\to Q_H$ for $N\subseteq H$ are those of
\cite[Proposition 3.14, Theorem 5.2 and Proposition 5.5]{A-DGShadows}.
We use the continuous isomorphism
\[
 \Psi:GT_0\xrightarrow{\ \sim\ }\varprojlim_{N\in\mathcal I}Q_N,
 \qquad (\lambda,f)\longmapsto
 \bigl([ (\lambda-1)/2,f]_N\bigr)_N.
\]
For a positive integer $d$, write $T_d(Q)=\{q\in Q:q^d=1\}$.
For each $H\in\mathcal I$, the image of $T_d(GT_0)$ in $Q_H$ is exactly
\[
 \bigcap_{\substack{N\in\mathcal I\\N\subseteq H}}
 R_{N,H}\bigl(T_d(Q_N)\bigr).                         \tag{D.23a}
\]
An element of order exactly $d$ in this intersection has a lift of
order exactly $d$. Every element of order exactly $d$ in $GT_0$ has
an image of order exactly $d$ at some isolated level.
For $a,b\in GT_0$, conjugacy in $GT_0$ is equivalent to conjugacy
of $\Psi(a)_N$ and $\Psi(b)_N$ in every $Q_N$.
\end{claim}
\begin{proof}
The cited theorem identifies the two profinite groups, including their
topologies. Its Proposition 5.1 makes the coordinates cofinal in the
finite quotients of $\widehat\Z$ and $\widehat F_2$. The map between
the normalization in that paper and ours is exactly
$(\widehat m,\widehat f)\mapsto(2\widehat m+1,\widehat f)$;
Proposition 5.5 checks the equations and the multiplication under this
map. Thus its inverse limit concerns the full $GT_0$ of Part E.

Fix $q\in Q_H$ in the right side of (D.23a). For every isolated
$N\subseteq H$ define the finite nonempty set
\[
 S_N=\{v\in Q_N:v^d=1,\ R_{N,H}(v)=q\}.
\]
Every reduction $R_{N,M}$ restricts to a function $S_N\to S_M$.
The inverse limit of these sets is nonempty. To verify the exact
compactness step, any finite list of coordinates and compatibility
conditions is dominated by one isolated subgroup
$L\subseteq H\cap N_1\cap\cdots\cap N_t$; Proposition 3.14 supplies
an isolated refinement of the intersection. Choose an element of
$S_L$ and project it to the specified coordinates. This satisfies
all the finitely many conditions. The closed conditions in the
compact product $\prod_NS_N$ consequently have the finite
intersection property. A compatible element maps under
$\Psi^{-1}$ to $g\in GT_0$ with $g^d=1$ and $g_H=q$.
The reverse inclusion follows by taking the coordinates of such a
$g$. The levels below $H$ are cofinal, so the constructed family is
a full inverse-limit point.

When $q$ has order $d$, the order of $g$ both divides $d$ and is a
multiple of the order of its image $q$, hence equals $d$.
Conversely, for $g$ of order $d$, choose, for each prime $p\mid d$,
a finite coordinate detecting $g^{d/p}\ne1$. A common isolated
refinement detects all these powers. The order of its image divides
$d$ and divides none of the proper divisors $d/p$, so it is $d$.

For the conjugacy assertion, the sets
\[
 C_N(a,b)=\{v\in Q_N:v\Psi(a)_Nv^{-1}=\Psi(b)_N\}
\]
are finite nonempty sets with restriction maps. The same finite
intersection argument gives a compatible conjugator, which
$\Psi^{-1}$ sends to a conjugator in $GT_0$. Necessity follows by
projection. The construction does not assume surjectivity of a
reduction map or that every element of $Q_N$ comes from $GT_0$.
\end{proof}

\begin{claim}{D.24}{Proved}{Dihedral tests with an order bound and the full dyadic quotient}
The program \texttt{code/part\_a/gt\_shadow\_refinement.py} executes
the finite order-two and order-three tests below. Each entry counts
nonidentity elements with order dividing $d$ at the lower level,
and then those admitting a lift with order dividing $d$ at the
upper level.
\begin{center}
\begin{tabular}{rrrrr}
\toprule
$n$&$q$&$d$&Lower elements&Elements admitting such a lift\\
\midrule
6&18&2&7&7\\
6&18&3&2&0\\
18&54&3&8&0\\
54&162&3&8&0\\
14&98&3&14&14\\
4&8&2&3&1\\
8&16&2&7&2\\
16&32&2&11&4\\
32&64&2&19&8\\
\bottomrule
\end{tabular}
\end{center}
The levels are isolated dihedral shadows $Q_{K^{(n)}}$. Their
subsystem with $n=2^\alpha$, $\alpha\geq2$, has inverse limit
\[
 H_2=\{(k,(-1)^a5^b)\in
 \Z_2\rtimes\Z_2^\times:
 a\in\{0,1\},\ b\in\Z_2,\ k\equiv b\pmod2\}.
                                                               \tag{D.24a}
\]
Every finite subgroup of $H_2$ is trivial or cyclic of order two,
and all its order-two subgroups are conjugate to
$\langle(0,-1)\rangle$. There is a continuous surjection
$GT_0\to H_2$. The family of all dihedral levels is not cofinal
in $\mathcal I$; an explicit $A_5$ quotient below proves this.
\end{claim}
\begin{proof}
Write $D_n=\langle r,s\mid r^n=s^2=1,\ srs^{-1}=r^{-1}\rangle$.
The homomorphism
\[
 \psi_n:PB_3=\langle x,y\rangle\times\langle c\rangle
       \longrightarrow D_n^3,
 \quad x\longmapsto(r,s,s),\quad
 y\longmapsto(rs,r,rs),\quad c\longmapsto1
\]
defines $K^{(n)}=\ker\psi_n$. For even $n\geq4$, the explicit
finite model of \cite[Theorem 4.3 and Proposition 4.5]{A-BDHPShadows}
is
\[
 Q_{K^{(n)}}\simeq
 \left\{(k,u)\in(\Z/(n/2))\rtimes(\Z/(2n))^\times:
 \begin{array}{l}
 k\equiv\nu(u)\pmod2\text{ when }4\mid n
 \end{array}\right\},                                  \tag{D.24b}
\]
where $\nu(u)=(u-1)/4$ for $u\equiv1\pmod4$ and
$\nu(u)=(u+1)/4$ for $u\equiv3\pmod4$.
Multiplication is $(k,u)(l,v)=(k+ul,uv)$, with the indicated
moduli. For $n\mid q$, reduction is
$(k\bmod(q/2),u\bmod2q)\mapsto(k\bmod(n/2),u\bmod2n)$.

The code constructs the actual subgroup
$G_n=\langle(r,s,s),(rs,r,rs)\rangle\leq D_n^3$ for
$n=4,6,8,12,14$. It computes its commutator subgroup and exhausts
all pairs $(m,f)$ with $0\leq m<n$, $\gcd(2m+1,n)=1$ and
$f\in G_n'$. It tests the two cycle equations
\[
 \theta(f)f=1,\qquad
 \tau^2(y^mf)\tau(y^mf)y^mf=1,
\]
where $\theta(x)=y$, $\theta(y)=x$, $\tau(x)=y$ and
$\tau(y)=(xy)^{-1}$. It checks descent of these maps, and of
$x\mapsto x^{2m+1}$, $y\mapsto f^{-1}y^{2m+1}f$, on every
labelled edge of the finite Cayley graph, and checks that the
resulting maps are bijective. The accepted pairs are exactly
those in (D.24b), of respective orders $4,12,16,24,84$.
This checks the reduction to the quoted finite model on the
tested groups, with no reliance on floating point arithmetic.
Projection of $G_6,G_{12},G_{14}$ onto $D_3,D_3,D_7$
exhibits finite quotients in which the odd-order rotation and
order-two reflection do not commute; their nilpotence would
force these coprime Sylow subgroups to commute. These tests
therefore include nonnilpotent finite quotients of $\widehat F_2$.

For the refinement counts, the code enumerates every element of
(D.24b), computes its $d$th power using the affine multiplication,
reduces every upper root, and counts the exact fibres. For example,
at $n=6$ the group is
$\operatorname{Aff}(\F_3)\times C_2$, and its two order-three
elements are the nonzero translations. At $q=18$, an element of
order dividing three has multiplier $u\equiv1,4,7\pmod9$ and
translation divisible by $3$. Both its affine coordinates reduce
to the identity modulo $3$. Thus neither order-three element
at $n=6$ has an order-dividing-three lift, although the unbounded
reduction map is surjective. At $n=14$ the multipliers $2,4$
modulo $7$ have order three. Each lifts to its order-three
root modulo $49$, and each of the seven lifts of a translation
coordinate is allowed because $1+u+u^2=0$ there. This gives
seven order-three lifts of each of the fourteen lower elements,
as the exhaustive output records. The same program gives the
remaining rows and the full table in its accompanying CSV.

The compatible affine coordinate reductions identify the dyadic
inverse limit with (D.24a): the residue of $b$ modulo $2$ is the
condition in (D.24b), and the inverse limits of the two coordinate
rings are $\Z_2$ and $\Z_2^\times$. The surjection from $GT_0$
is \cite[Theorem 5.7 and Proposition 5.8]{A-BDHPShadows}; more
precisely, that theorem proves surjectivity already after
restriction along $G_{\Q}\to GT_0$.

The translation kernel $\Z_2$ is torsion-free, so a finite
subgroup of $H_2$ injects into the finite-order subgroup of
$\Z_2^\times$. This subgroup is $\{1,-1\}$: the parametrization
$\Z_2^\times=\{\pm1\}\times5^{\Z_2}$ has torsion-free second
factor, as follows from
$v_2(5^{2^j}-1)=j+2$ and passage to the inverse limit.
Every nonidentity finite-order element of $H_2$ is consequently
$(k,-1)$ with $k\in2\Z_2$. Put $t=k/2$ and choose
$v=5^{t\bmod2}$. Then $(t,v)\in H_2$, and direct multiplication
gives
\[
 (t,v)(0,-1)(t,v)^{-1}=(2t,-1)=(k,-1).
\]
This proves the full finite-subgroup and conjugacy assertions
for this infinite quotient, rather than extrapolating the
finite dyadic counts.

Finally every $G_n$ is metabelian: the commutator subgroup of
$D_n^3$ lies in its abelian rotation subgroup, and $G_n'$ is
a subgroup of that rotation subgroup. Define the epimorphism
$\eta:F_2\to A_5$ by
$x\mapsto(1\ 2\ 3\ 4\ 5)$ and $y\mapsto(1\ 2\ 3)$.
The code enumerates its image and derived subgroup; both have
order $60$, proving directly that this image is a nontrivial
perfect group. Proposition 5.1 of \cite{A-DGShadows} gives an
isolated $N$ with $N\cap F_2\subseteq\ker\eta$.
An inclusion $K^{(n)}\subseteq N$ would give the surjection
\[
 G_n=F_2/(K^{(n)}\cap F_2)
   \longrightarrow F_2/\ker\eta=A_5.
\]
Its image would be metabelian, contradicting the computed
nontrivial perfect derived subgroup. Hence no such inclusion
exists. This proves the asserted failure of cofinality by an
explicit quotient and the exact factoring map.
\end{proof}

These tests give explicit finite obstructions for D.12.
D.32 proves, by its full fixed-subgroup and finite-cover argument,
that $GT_0$ has no nonidentity element of odd order and that every
finite subgroup has order at most two, detected by the dyadic
multiplier. D.34 further proves that a nonidentity involution
has the global multiplier $-1$.
The fourteen order-three shadows at $K^{(14)}$ survive the
executed test through $K^{(98)}$. Combining D.32 with D.23
shows that each fails the order-three lifting condition at some
isolated refinement; a common refinement excludes all fourteen.
This existence conclusion does not identify that refinement or
alter the recorded finite counts. D.25 explains why the entire
dihedral subsystem alone cannot supply the exclusion.
The remaining full $GT_0$ question is conjugacy of its
involutions. D.35 gives the exact third-cusp equation and the
finite conjugator criterion, and D.44 supplies a coboundary
satisfying its first symmetry. D.42 also proves the global
multiplier restriction for every finite subgroup of the full
three-cusp loop-outer group.
